% Descripcion del algoritmo
\section{Algorithm}
\label{sec:algorithm}

The algorithm loads the ULDs one at a time. Each ULD is loaded by BSG, a beam search guided by a
greedy algorithm, which builds the load with blocks of boxes placed in maximal free spaces
(Section~\ref{sec:alg-bsg}). The ULD-specific features (the bevelled corners, the edge and the
substructure, and the support and fragility constraints) are handled in the generation and the
placement of the blocks (Section~\ref{sec:alg-uld}), and the CoG constraint at several levels
of the search (Section~\ref{sec:alg-cog}). Section~\ref{sec:alg-multi} describes the
multi-ULD framework.

\subsection{Loading one ULD with BSG}
\label{sec:alg-bsg}

\paragraph{Blocks.}
A block is a cuboid arrangement of boxes that is placed as a unit. Before loading a ULD, a set
of blocks is generated from the boxes still available, as in \citet{Fanslau2010} and \citet{Araya2014}: the blocks of
one box are the boxes in their allowed orientations, and new blocks are obtained by joining two
existing blocks along the $x$, $y$ or $z$ axis, so that every block has a guillotine structure.
A joined block is kept only if its boxes are available, it fits in the bounding box of the ULD,
the boxes fill at least a proportion $f_{\min}$ of its volume, and it is different from the
blocks already generated (same dimensions and same boxes); the generation stops when
$N_{\max}$ blocks have been generated. We use $f_{\min} = 0.98$, as in BSG
\citep{Araya2014}, so that the blocks are almost solid, and $N_{\max} = 500$
(Section~\ref{sec:exp-setup}). Each block stores its weight and its CoG, computed recursively from the two
blocks that form it, which allows us to evaluate the CoG of a load without visiting its boxes.

\paragraph{Free spaces and actions.}
The free space of the ULD is represented by the set of maximal free spaces, the largest
cuboids that do not intersect the placed blocks \citep{Parreno2008}. When a block is placed, the
free spaces that intersect it are replaced by their maximal parts outside the block, and the
spaces too small for any remaining box are discarded. In \citet{Araya2014}, each space is anchored at its corner
closest to a corner of the container, and the next space to fill is the one whose anchor corner
is closest (in Manhattan distance) to the corresponding corner of the container, so that the
load grows from all the corners towards the center. Here, instead, every space is anchored at
its corner closest to the origin of the ULD and the next space is the one closest to the
origin, so that the load grows from one corner and stays grouped: a grouped load leaves room to
translate it, which is needed to center its CoG (Section~\ref{sec:alg-cog}), and when it does
not fill the ULD it is not spread against several walls. Ties are broken by the larger volume.
An action places a block that fits in the selected space at its anchor corner, which may be
modified as explained in Section~\ref{sec:alg-uld}; a block is a candidate only if its weight
does not exceed the remaining weight capacity of the ULD. A state of the
search is a partial loading of the ULD, with its free spaces and the remaining boxes.

\paragraph{Evaluation of the blocks.}
The blocks that fit in the selected space are evaluated with the VCS function of
\citet{Araya2017}:
\begin{equation}
\label{eq:vcs}
\mathrm{VCS}(b) = V(b)\,\bigl(1 - \mathrm{loss}(b)\bigr)^{\beta}\,\mathrm{CS}_p(b)^{\alpha}\,
	n(b)^{-\gamma},
\end{equation}
where $V(b)$ is the volume of the boxes of the block, $\mathrm{loss}(b)$ is the proportion of
the volume of the free space that cannot be filled with the remaining boxes after placing the
block (estimated with one-dimensional knapsack problems on each dimension of the space),
$\mathrm{CS}_p(b)$ is the proportion of the surface of the block in contact with the placed
blocks and the walls of the ULD, where faces at a distance of at most a proportion $p$ of the
block dimension are considered in contact, and $n(b)$ is the number of boxes of the block.

\paragraph{Greedy and beam search.}
The greedy algorithm repeatedly selects the next free space and places the best evaluated block
in it, until no block can be placed. The value of a complete loading is its volume utilization
penalized by the CoG excess (Section~\ref{sec:alg-cog}). BSG \citep{Araya2014} is a beam
search in which the children of each node are generated with the best evaluated blocks for the
selected space, and each child is evaluated by the value of its greedy completion. At each
level, the $w$ children with the best greedy values are kept, and a total of at most $m$
children are generated, so that each node of the beam is expanded with about $m/w$ blocks. The
best complete solution found by any greedy completion is returned. BSG is run iteratively with
increasing effort until the time limit $t$ is reached \citep{Araya2017}: it starts with $w = 4$
and $m = 16$, and after each run $w$ and $m/w$ are multiplied by $\sqrt{2}$. The search also stops
as soon as a solution loads all the available boxes with the CoG in range, since no better
solution exists for that ULD. The implementation follows
the optimizations proposed by \citet{Guesser2026} for BSG: the free spaces and the placed
blocks are stored in contiguous arrays, the contact surface is computed only with the blocks
near the selected space, the duplicate blocks are detected with a hash during the generation,
and the greedy completions are memoized: the states visited by the complete greedy runs are
stored (identified by a hash of their placed blocks), and when a greedy run reaches a known
state, the rest of the solution is rebuilt from memory instead of evaluating the blocks again.
In addition, the volume loss of a block is computed only when an upper bound of its evaluation
can enter the best ones. None of these optimizations changes the result of the search.

\subsection{ULD constraints}
\label{sec:alg-uld}

\paragraph{Bevelled corners.}
Each cut of a ULD is a half-plane $a_j x + c_j z \le d_j$ in the $x$-$z$ plane, and the usable space
is the intersection of the bounding box $[0, L] \times [0, W] \times [0, H]$ with all the cuts. A
block $[x, x+\ell] \times [z, z+h]$ (in the $x$-$z$ plane) satisfies cut $j$ if and only if its
corner that maximizes $a_j x + c_j z$ does, that is, if
$a_j (x + \ell\,[a_j > 0]) + c_j (z + h\,[c_j > 0]) \le d_j$. A block placed at the anchor corner
of a space may cross one of the cuts. In that case, the block is shifted along the $x$ axis,
inside the space, to the closest position that satisfies all the cuts: for a fixed $z$, each
cut with $a_j > 0$ gives an upper bound and each cut with $a_j < 0$ a lower bound on $x$,
\[
x \le \Bigl\lfloor \frac{d_j - c_j (z + h\,[c_j > 0])}{a_j} \Bigr\rfloor - \ell \quad (a_j > 0),
\qquad
x \ge \Bigl\lceil \frac{d_j - c_j (z + h\,[c_j > 0])}{a_j} \Bigr\rceil \quad (a_j < 0),
\]
and the block is moved to the point of the resulting interval, intersected with the extent of
the space, closest to the anchor; if the interval is empty, the block is not placed in the
space. Since the cuts are
orthogonal to the $x$-$z$ plane, shifting the block along $x$ keeps it on the same base, which
is necessary for its support. The free spaces are still cuboids, so a space next to a cut may
contain positions that are not usable; they are simply not used by any block.

\paragraph{Edge and substructure.}
With a floor edge of width $e$ and offset $\delta$, the initial free spaces are the space above
the floor inside the edge, $[e, L-e] \times [e, W-e] \times [0, H]$, and the space above the height
$\delta$ over the whole floor, $[0, L] \times [0, W] \times [\delta, H]$. With a substructure, the
only initial space is the second one, and the floor inside the edge is raised to $\delta$. A space above the edge has no support until boxes have been loaded inside
it; instead of discarding it, it is kept aside (\emph{dormant}) and reconsidered after the next
block is placed. Each ULD with an edge is loaded twice, on its floor and on the substructure,
and the best of the two loadings is kept; the blocks (single boxes and composite blocks) are
generated for each of the two loadings, since the space in which they must fit is different.

\paragraph{Support and fragility.}
The support criterion of the problem (the four corners of the base, or the area criterion of
Section~\ref{sec:problem}) is checked box by box. Let $R$ be the base of a box placed at height
$z > 0$, and let the \emph{direct} faces be the top faces of the placed boxes at height $z$ and the
\emph{padding} faces those at a height in $[z - \hbar, z)$. The box is supported if each of the
four corners of $R$ lies on a direct face of a non-fragile box or, with the area criterion, if
it overlaps some direct face and $A_d + A_p \ge o\,|R|$, where $A_d$ and $A_p$ are the areas of
$R$ covered by direct and padding faces of non-fragile boxes and $o$ is the minimum supported
proportion ($o = 0.9$ in the adapted settings); the raised floor of the substructure is also a
direct face. The box is infeasible if $R$ overlaps a direct face of a fragile box, and the
padding faces of fragile boxes give no support. Inside a block, the support is guaranteed by
construction: two blocks are joined along $z$ only if the lower block has a complete top face
without fragile boxes and the upper block fits on it, so that every box of the upper block has
its whole base, and in particular its four corners, on non-fragile boxes. Therefore, when a
block is placed on the base of the selected space, only the boxes at its bottom are checked. Since the corner criterion allows gaps under a box, we also forbid placing a fragile box in
such a gap if it would touch the base of the box above. If no block can be placed in a space
because of the support, the space is discarded or, with an edge, kept dormant.

\subsection{Center of gravity}
\label{sec:alg-cog}

The CoG constraint is taken into account in two ways, besides the anchoring of the spaces at
the origin.

\paragraph{Translation of the load.}
Once a ULD has been loaded, the whole load can be translated horizontally without changing its
feasibility, as long as every block stays inside the ULD and its cuts and, if the ULD has an
edge, the blocks below the height $\delta$ stay inside the edge and the blocks on the
substructure stay on it. These conditions define an interval of feasible translations in $x$
and in $y$: a block $[x_1, x_2] \times [y_1, y_2] \times [z_1, z_2]$ allows
$\Delta x \in [-x_1, L - x_2]$ and $\Delta y \in [-y_1, W - y_2]$, intersected, for $\Delta x$, with
the bounds that the cuts give (as above, with $z$ fixed) and, if $z_1 < \delta$ or the block lies
on the substructure, with $[e - x_1, L - e - x_2]$ and $[e - y_1, W - e - y_2]$. The intervals of
the load are the intersections over its blocks, so they are updated in constant time when a
block is placed. The final load of each ULD is translated by
$\Delta x = \min\bigl(\max(L/2 - \bar x, \Delta x^{\min}), \Delta x^{\max}\bigr)$, and $\Delta y$ in the
same way, the feasible amount that brings its CoG $(\bar x, \bar y)$ closest to the center of
the ULD.

\paragraph{CoG excess.}
Let $(\bar x, \bar y, \bar z)$ be the CoG of a loading after the best feasible translation
$(\Delta x, \Delta y)$. Its CoG excess is
\begin{equation}
\label{eq:excess}
\varepsilon = \frac{\bigl(|\bar x + \Delta x - L/2| - d^x\bigr)^+}{d^x}
	+ \frac{\bigl(|\bar y + \Delta y - W/2| - d^y\bigr)^+}{d^y}
	+ \frac{(\bar z - d^z)^+}{d^z},
\end{equation}
where $(a)^+ = \max(a, 0)$. The CoG constraint is satisfied if and only if $\varepsilon = 0$.

\paragraph{Value of a solution.}
The value of a (partial or complete) loading with volume utilization $u$ is
\begin{equation}
\label{eq:value}
v = \frac{u}{1 + \lambda\,\varepsilon},
\end{equation}
where $\lambda \ge 0$ is a parameter. This value is used to compare the greedy completions in
BSG and to select the best loading of a ULD. A multiplicative penalty keeps the value of any
non-empty loading positive, so that an unbalanced loading is still preferred to an empty ULD.
With a large $\lambda$, the CoG constraint becomes almost a hard constraint; with
$\lambda = 0$, it is ignored.

The weight and the moments of the partial loading and its interval of feasible translations
are stored in each state, so that the CoG excess of a greedy completion is computed in constant
time per placed block. Penalizing the CoG excess also in the evaluation of each candidate block
had no consistent effect in our experiments, and it is not used.

Algorithm~\ref{alg:load} summarizes how a ULD of type $k$ is loaded with the available boxes
$B$. If the ULD has an edge, two loadings are built, one on the floor and one on the substructure
(line~\ref{a1:for}). For each of them, the blocks that fit in that space are generated
(line~\ref{a1:gen}), BSG builds a loading within the time limit $t$ (line~\ref{a1:bsg}), the
loading is translated to center its CoG (line~\ref{a1:tr}), and the one with the larger value
$v$ (Equation~\ref{eq:value}) is kept (line~\ref{a1:best}).

\begin{algorithm}[htbp]
\caption{\textsc{LoadULD}$(k, B)$: load a ULD of type $k$ with the available boxes $B$.}
\label{alg:load}
\begin{algorithmic}[1]
\State $S^* \gets$ empty loading
\For{$s \in \{\text{floor}, \text{substructure}\}$} \label{a1:for}
	\Comment{the substructure only if the ULD has an edge}
	\State generate the blocks of $B$ that fit in a ULD of type $k$ loaded on $s$ \label{a1:gen}
	\State $S \gets \textsc{BSG}(k, s, B, t)$ \label{a1:bsg}
	\State translate $S$ to center its CoG \label{a1:tr}
	\If{$v(S) > v(S^*)$} $S^* \gets S$ \EndIf \label{a1:best}
\EndFor
\State \Return $S^*$
\end{algorithmic}
\end{algorithm}

\subsection{Loading several ULDs}
\label{sec:alg-multi}

Algorithm~\ref{alg:multi} outlines the multi-ULD framework. It has a sequential construction
phase, the insertion of the boxes left unloaded when the number of ULDs is limited, and a phase
that reduces the total volume of the ULDs, the objective of the problem. The construction
(lines~\ref{a2:cons}--\ref{a2:add}) loads one ULD at a time, choosing its type among the
candidates of the rule of \citet{Hessler2025} (line~\ref{a2:cand}) and, when the ULDs are
limited, loading first the boxes that would otherwise be left out (line~\ref{a2:prio}). The
boxes that remain are inserted into the loaded ULDs (line~\ref{a2:ins}); if some are still
unloaded, the construction is repeated with them in the set $P$ of boxes to load first, and the
solution with the least unloaded volume is kept (lines~\ref{a2:keep}--\ref{a2:restart}).
Finally, the reduction phase (lines~\ref{a2:red}--\ref{a2:undo}) tries to replace each ULD by
inserting its boxes into the others (line~\ref{a2:rins}) and loading the rest into smaller ULDs
(line~\ref{a2:new}). The three phases are detailed in the following paragraphs.

\paragraph{Sequential construction.}
The ULDs are loaded one at a time with \textsc{LoadULD}, and the type of the next ULD is selected
as in \citet{Hessler2025}: among the ULD types into which the boxes that fit in the fewest
available types fit, the type with the largest total volume of the remaining boxes that fit in
it is selected (ties are broken by the larger ULD volume). When this rule leaves several
candidate types (those that contain the boxes with the fewest options), each candidate is
loaded with the shorter time limit of the improvement phases (see below), the one that
maximizes $L - \theta V$ is chosen and it is loaded again with the time limit $t$, where $L$ is the loaded volume and $V$
the volume of the ULD, unless it leaves the CoG out of range and another candidate does not.
With $\theta = 0$ this is the type that loads the most volume; a small $\theta > 0$ prefers a
smaller ULD when it loads almost the same volume. A box fits in a ULD type if its
weight does not exceed the capacity and, in some allowed orientation, it can be placed alone on
the floor inside the edge or on the substructure. When the number of ULDs is limited, some boxes
may only fit in the selected type among the ULDs still available. If the loading of the ULD
leaves some of them out, the ULD is loaded again with these boxes first, and the load is then
completed with the other boxes, keeping the first ones at their positions. If, after the
insertion of the unloaded boxes into the loaded ULDs (see below), some boxes are still left
out, the construction is restarted, and the boxes left out in the previous constructions are
also loaded first, into the first ULD in which they fit. This handles boxes that fit in several
types but conflict with each other, for instance two large boxes that fit only in the largest
types and cannot share a ULD. The solution with the smallest unloaded volume is kept, with at
most three restarts.

\paragraph{Insertion of boxes into the loaded ULDs.}
Both phases after the construction rely on the same procedure, \textsc{Insert}, which inserts a pool of
boxes into the ULDs already loaded, without adding ULDs, by ruin and recreate
\citep{Schrimpf2000}. For $\rho = 0.3$, $0.6$ and $0.9$, in this order, a proportion $\rho$ of the
loaded volume of each ULD is removed, taking the boxes with nothing on top of them from the
highest one down, so that the remaining load is still feasible; then the ULDs, in decreasing
order of free volume, are completed with \textsc{LoadULD} with the removed boxes and the pool,
keeping the remaining boxes at their positions. A ULD whose new loading has the CoG out of range
when the old one had not keeps its old loading. The new solution is accepted if the volume of the
pool decreases, and the procedure stops as soon as the pool is empty. Reloading each ULD from
scratch with its boxes and the pool, or giving priority to the boxes of the pool, did not improve
the results at the same running time.

\paragraph{Unloaded boxes.}
When the number of ULDs is limited and some boxes remain unloaded after the construction,
\textsc{Insert} is applied to them.

\paragraph{Reduction of the volume.}
The construction tends to use the largest ULDs, and its last ULDs are often far from full. The
reduction phase replaces ULDs by cheaper ones with a single move that generalizes the elimination
of ULDs. The ULDs are considered in decreasing order of lost volume (the volume of the ULD minus
the volume of its boxes). The boxes of the selected ULD $S$ are freed and inserted into the other
ULDs with \textsc{Insert}. If all of them are inserted, $S$ is eliminated. Otherwise, the remaining
boxes are loaded into one or two new ULDs whose total volume is smaller than that of $S$: the
combinations of one or two available types with enough volume and weight capacity, in which every
remaining box fits, are tried in increasing order of total volume, loading the largest ULD first
with all the remaining boxes and the second one with the rest. The first combination that loads
all the boxes, without leaving a CoG out of range when the CoG of $S$ was not, replaces $S$;
otherwise the changes made by \textsc{Insert} are undone and the next ULD is tried. After every
change the phase starts again, and it stops when no ULD can be changed. Every change reduces the
total volume of the ULDs, so the phase terminates; a change may increase the number of ULDs, when
one ULD is replaced by two smaller ones. Since \textsc{Insert} and the reduction load many ULDs
again, BSG is given a shorter time limit in these phases than in the construction ($0.1\,t$ in
our experiments).

\begin{algorithm}[htbp]
\caption{Multi-ULD framework.}
\label{alg:multi}
\begin{algorithmic}[1]
\State $P \gets \emptyset$ \Comment{boxes to load first (restarts)}
\Repeat
	\State $B \gets I$; $\mathcal{S} \gets \emptyset$ \Comment{remaining boxes, loaded ULDs}
	\While{$B \neq \emptyset$ \textbf{and} some available ULD type fits a box of $B$} \label{a2:cons}
		\State $K \gets$ candidate ULD types for $B$ \citep{Hessler2025} \label{a2:cand}
		\State $S \gets$ the loading $\textsc{LoadULD}(k, B)$, $k \in K$, that maximizes $L - \theta V$
		\If{the ULDs are limited and $S$ leaves out boxes of $P$ or that only fit in its type} \label{a2:prio}
			\State load these boxes first and complete $S$ with the other boxes of $B$
		\EndIf
		\State $\mathcal{S} \gets \mathcal{S} \cup \{S\}$; remove the boxes of $S$ from $B$ \label{a2:add}
	\EndWhile
	\If{$B \neq \emptyset$}
		\State $\textsc{Insert}(\mathcal{S}, B)$ \Comment{unloaded boxes (limited ULDs)} \label{a2:ins}
	\EndIf
	\State keep $(\mathcal{S}, B)$ if it is the first or leaves less volume unloaded; $P \gets P \cup B$ \label{a2:keep}
\Until{$B = \emptyset$ or three restarts} \label{a2:restart}
\State $(\mathcal{S}, B) \gets$ the kept solution
\Repeat \Comment{reduction of the volume} \label{a2:red}
	\For{$S \in \mathcal{S}$ in decreasing order of lost volume}
		\State $B' \gets$ boxes of $S$; $\textsc{Insert}(\mathcal{S} \setminus \{S\}, B')$ \label{a2:rins}
		\If{$B' = \emptyset$} remove $S$; \textbf{break} \EndIf
		\If{$B'$ fits in 1 or 2 new ULDs of smaller total volume than $S$} \label{a2:new}
			\State replace $S$ by them; \textbf{break}
		\EndIf
		\State undo the changes of \textsc{Insert} \label{a2:undo}
	\EndFor
\Until{no ULD was changed}
\State \Return $\mathcal{S}$, $B$
\end{algorithmic}
\end{algorithm}
